\subsection{What is established, and what is open}
\label{sec:open}

\paragraph{Established.} Full-state storage fails by an exact count (Section~\ref{sec:walls}). The cut entropies of the evolved state are measured for $N=2$--\num{Nmax_measured}, the pair-counting law fails at $N=5,6$, and the correlation graph reproduces the entropies (Section~\ref{sec:measured}). For $N\le5$ the start is two eigenstates plus a dressing at $1$--$3J$, an exact reflection symmetry removes the other levels, and a control run shows that most of the cut entropy is present without any transfer (Section~\ref{sec:static}). If the start stays in the span of two states and $M$ sidebands, the signal is a short sum of static cosines and the Schmidt rank is bounded (Proposition~\ref{prop:cond}).

\paragraph{Not supported.} A classical--quantum window for the \emph{dynamics}: the tensor-network wall prices the entanglement of the evolved state, most of which is static, and the quantum edge is a leading-order price of one evolution (without preparation, compilation overheads or repeated measurements) for a window $\num{win_ratio}$ times shorter than the one in which the entropies were maximized. The walls at $N\approx\num{wall_sat_grid}$ rest on a scenario for the one-particle entropy beyond $N=\num{Nmax_measured}$, and the sharing law $w_{ij}\approx2\bar S_1/(N-1)$ describes $N\le6$ but is not derived.

\paragraph{Open.}
\begin{itemize}[itemsep=0.15em]
\item \emph{The dressed-transfer spectrum} decides the cost of the task: the number $M(N)$ of sidebands and the cost of their energies, weights and matrix elements (Section~\ref{sec:openproblem}). A Gaussianity (Huang--Rhys) test of the pair weight at $k\ge16$, which needs no propagation, and a converged spectrum at $N=6,7$ would settle most of it; neither was run.
\item \emph{Entropy beyond $N=\num{Nmax_measured}$.} The $N=6$ point rules out the pair law but does not separate the sharing-law scenarios (Table~\ref{tab:scen}); a direct matrix-product-state evolution at $N=6$--$10$ would fix the growth. A bound relating the rank of a truncated representation to the entropy, for the observable rather than the full state, is also missing.
\item \emph{Imaginary-time Monte Carlo} is the nearest classical rival (Section~\ref{sec:mc}). It is the right picture at $N=2,3$ (Table~\ref{tab:mc}); at $N=4,5$ the signal needs several cosines, and whether this stays a few-cosine problem for $N\gtrsim10$, and whether the regime itself continues that far (Fig.~\ref{fig:tilt}a), decides whether anything is left of the window.
\item \emph{Local-information algorithms} that evolve the observable in the Pauli basis~\cite{begusic2024,begusic2025,rudolph2025,schuster2025} follow dynamics that look hard by entanglement; the reactivity of Ref.~\cite{schuster2026reactivity} separates cheap from expensive cases, and given Section~\ref{sec:static} a low value is now the likelier outcome for $N\le5$. Nothing here tests this prediction, and localization~\cite{nandkishore2015} is not suggested by a clean model with a single bias.
\item \emph{Hardness} can only be a property of a family, and for the fixed family here the observable can be tabulated (Section~\ref{sec:openproblem}). Quantum walks and interacting particles on graphs are universal~\cite{childs2009walk,childs2013walk}, but we have not shown that the capped two-body double-well family embeds such a construction, and a hardness theorem would not place a wall at a given $N$ without a conjecture of the form ``no classical algorithm beats $2^{cN}$''.
\end{itemize}

\paragraph{Limitations.} Entropies fall with $k$ and Schmidt ranks rise; neither is converged, and the pair weight converges at $N=4$ but not at $N=5$ (Table~\ref{tab:k}, Fig.~\ref{fig:dressing}a). All entropies are maxima over a window of about $\num{win_ratio}$ transfer times, whereas the quantum count prices one evolution; the shorter window has lower entropy (Table~\ref{tab:protocol}), so the comparison charges the classical side for the longer window, and a maximum over a finite window is a lower bound. The Lennard-Jones core is marginally resolved on the grid, so splittings and the resonance tilt depend on it (Table~\ref{tab:params}). The correlated start is the conservative definition for the classical side and its preparation is not priced for the quantum side. Particles are distinguishable and two-dimensional with one $\epsilon$ and one $\sigma$; identical bosons and three dimensions change numbers in ways we do not know. Spectra of the start exist for $N\le5$ (and $N=6$ in one unconverged basis); the closed two-state structure is measured at $N=3,4$ and, at $N=5$, only the weights.

\FloatBarrier
\section{Conclusion}
\label{sec:conclusion}

We asked when a quantum computer is needed for real-time tunneling in a small cluster and answered with inequalities whose inputs are measured where they can be. Exact storage fails by a count (imaginary-time Monte Carlo is not excluded). For tensor networks the cost is set by the entropy of the cut that every layout must pay; we measured it up to $N=\num{Nmax_measured}$, found that it grows by less per added particle than pair counting predicts, and described it with a correlation graph. Carried to larger $N$ under stated scenarios this places a tensor-network wall between $N=\num{wall_pess_grid}$ and $N=\num{wall_opt_grid}$ on the grid, and a leading-order count prices one evolution on a quantum computer at a mass edge that falls slowly with $N$.

The spectrum of the start changes what these numbers mean. For $N\le5$ the prepared state is two eigenstates plus a dressing by soft internal levels, an exact symmetry removes the other levels, and most of the cut entropy is present without any tunneling. The tensor-network walls therefore price the static entanglement of the isomers. If the structure persists, the signal is a short sum of cosines of static quantities, and the cost of the task is the cost of that spectrum. We have not shown that the structure persists, nor that it fails; the dressing grows with $N$ at fixed basis, the number of cosines needed grows with $N$, and the tunneling regime itself ends near $N\approx\num{fu_Ntop_lo}$--\num{fu_Ntop_hi} at fixed trap depth. A classical--quantum window is therefore not established by this report, and the question it was meant to answer, whether real-time tunneling in a cluster is classically hard, is open. What would settle it is a Gaussianity test of the dressing and a converged spectrum at $N=6,7$.

